\section{Ensemble implementation}
\subsection{Numerical fitting}
\label{app:numerics}
The Selected model, Inverse error mixture, and Fitted mixture use the same
$K$ input and fine field pairs. Their definitions appear in
Section~\ref{sec:method}. For joint fitting, the empirical error matrix is
symmetrized and divided by
$\max(\operatorname{tr}(\widehat C)/M,10^{-30})$. Here the trace sums diagonal
errors and $M=9$ counts models. This positive scaling changes the objective
magnitude without changing the matrix condition number or the exact minimizer. SLSQP uses the analytic
gradient, bounds $0\leq w_m\leq1$, the unit sum constraint, tolerance
$10^{-12}$, and at most 300 iterations. Uniform weights and the lowest
diagonal singleton initialize separate solves. Finite solutions with unit sum
residual below $10^{-6}$ are clipped to nonnegative values and renormalized.
Feasible starting points remain candidates, and the lowest fitting objective
determines the retained weights. These checks establish numerical feasibility,
not generalization to new inputs.

Uniform averaging, used as a control in Table~\ref{tab:main}, assigns every
member weight $1/M$. Forming the error matrices costs $O(KM^2P)$ for $P$
values per field. Streaming predictions and keeping only these matrices
requires $O(KM^2)$ storage. These costs exclude surrogate training and inference.
The companion package preserves additional regularization, pair fitting, and
coarse input refresh experiments separately from the three main rules.

\subsection{From a trained surrogate library to a deployed ensemble}\label{app:algorithm}
\refstepcounter{algorithm}\label{alg:calibration}
\begin{center}
\fbox{\begin{minipage}{0.94\linewidth}\small
\textbf{Algorithm \thealgorithm: train models, fit the ensemble, and predict}\par
\textbf{Input:} LF/HF training tables, $K$ reserved fine cases, fixed model recipes,
and one prescribed rule: Selected model, Inverse error mixture, or Fitted mixture.\par
1. Check input identities and working grids within the dataset. Train the surrogate models
using training data only. Record unresolved units or simulation metadata.\par
2. Predict fields for the fitting inputs and compute $G_1,\ldots,G_K$.\par
3. Fit the prescribed rule on all $K$ cases using the definitions in Section~\ref{sec:method}.\par
4. Save preprocessing, model identities, and fitted weights.\par
\textbf{Predict:} combine model predictions at new parameters using Equation~\eqref{eq:mixture}.\par
\textbf{Evaluate:} score reserved fine answers only after saving these choices.
\end{minipage}}
\end{center}
The algorithm separates surrogate training from ensemble fitting. The bundled
program implements steps 2 to 4 and combines saved model predictions for new
inputs. Dataset adapters and base model training remain separate components.
The three rules use the same fitting and evaluation cases within each comparison.


